% !TEX root = ../../Extensions of IQC Duality to Partial Integral Equations.tex

This section connects the duality framework developed above with
the established theory of finite-dimensional dynamic IQC multipliers. For
rational multipliers satisfying the strict positive--negative (PN)
conditions, Seiler~\cite{6915700} showed that a $J$-spectral factorization
provides a stable and stably invertible hard factorization, thereby converting
the associated soft, infinite-horizon IQC into a time-domain constraint that
holds on every finite horizon. As discussed in the linear parameter-varying (LPV) setting by Venkataraman and
Seiler~\cite{venkataraman}, this $J$-spectral factorization satisfies the filter
assumptions needed to apply the duality theory. The first part of this section
reviews this connection between strict PN multipliers, hard factorizations,
and duality. The second part illustrates the construction using a delayed
firing-rate model.
\subsection{Positive--Negative Multipliers}

Using the terminology of Definition~\ref{def:factorization}, we denote the
$J$-spectral factor by $\Theta$. The following result restates
\cite[Thm.~4]{6915700} in this notation. Moreover, for finite dimensional mutlipliers we can
define $\Pi = \Psi^\sim V\Psi\in\mathbb{RL}_\infty$, where $\Psi(j\omega)^\sim = \Psi(-j\omega)^\top$.
\begin{theorem}\label{thm:strict-pn-j-spectral}
Let $\Pi=\Pi^\sim\in\mathbb{RL}_\infty^{(n_{{z}}+n_{{w}})\times(n_{{z}}+n_{{w}})}$
where $\Pi_{11}\in\mathbb{RL}_\infty^{n_{{z}}\times n_{{z}}}$ and
$\Pi_{22}\in\mathbb{RL}_\infty^{n_{{w}}\times n_{w}}$. If $\Pi_{11}(j\omega)\succ0$,\,$\Pi_{22}(j\omega)\prec0$,
for all $\omega\in\mathbb{R}\cup\{\infty\}$,
then
\begin{enumerate}
    \item $\Pi$ admits a $J_{n_z,n_w}$-spectral factorization
          $\Pi=\Theta^\sim J_{n_z,n_w}\Theta$, where
          $\Theta,\Theta^{-1}\in
          \mathbb{RH}_\infty^{(n_z+n_w)\times(n_z+n_w)}$;
    \item $(\Theta,J_{n_z,n_w})$ is a hard factorization of $\Pi$.
\end{enumerate}
\end{theorem}
A Riccati-based construction of $J$-spectral factors for strict PN
multipliers is given in~\cite[App.~B]{Pfifer2015}, using rational
$J$-spectral factorization results~\cite{MEINSMA1995243}.

The distinction between a frequency-domain factorization identity and the
finite-horizon conditions used in Definition~\ref{def:factorization} is relevant
here. Preservation of both graph and inverse-graph inequalities is studied
through doubly-hard factorizations~\cite{carrasco2018equivalence}.
Finite-horizon IQCs with terminal costs provide another connection to
dissipativity~\cite{scherer2018dynamicdissipation}, while multiplier storage
can reduce conservatism~\cite{magicalhiddenenergy,hu}.
These results provide context for the finite-horizon assumptions imposed in
Definition~\ref{def:factorization}; they do not by themselves establish the
truncated quadratic-form identity in that definition.


\subsection{Example: Generation of $\beta$-Oscillations in the Subthalamic Nucleus-Globus Pallidus Network}
Consider the delayed firing-rate model in
\cite[Eq.~(5)]{nevadoHolgado2010betaOscillations}, and define the stimulation
input $\tilde u(t):=I_{\mrm e}\mrm{Stim}(t)$.
\begin{equation}
      \begin{aligned}
      \tau_{\mrm{S}} \dot{\mrm{STN}}(t)
          ={}&F_{\mrm{S}}\!\bigl(-w_{\mrm{GS}}\mrm{GP}
              (t-\tau_{\mrm{GS}})\\
             &\qquad\quad +w_{\mrm{CS}}\mrm{Ctx}\bigr)
              -\mrm{STN}(t)+\tilde u(t),\\
      \tau_{\mrm{G}} \dot{\mrm{GP}}(t)
          ={}&F_{\mrm{G}}\!\bigl(w_{\mrm{SG}}\mrm{STN}
              (t-\tau_{\mrm{SG}})\\
             &\qquad\quad-w_{\mrm{GG}}\mrm{GP}(t-\tau_{\mrm{GG}})\\
             &\qquad\quad-w_{\mrm{XG}}\mrm{Str}\bigr)-\mrm{GP}(t).
      \end{aligned}
\end{equation}
Let $\mrm{STN}$ and
$\mrm{GP}$ denote the physical firing rates of the SubThalamic Nucleus
(STN) and Globus Pallidus pars Externa (GPe). The activation functions are
\begin{equation}\label{eq:stn-gpe-uncentered-sigmoid}
F_i(q)=\frac{M_i}{1+\exp(-4q/M_i)(M_i-B_i)/B_i},
\qquad i\in\{S,G\},
\end{equation}
so $F_i(0)=B_i$. Define the sigmoid midpoint and its centered form as
\begin{equation}\label{eq:stn-gpe-centered-sigmoid}
\begin{aligned}
c_i&:=\frac{M_i}{4}\ln\!\left(\frac{M_i-B_i}{B_i}\right),\\
\sigma_i(q)&:=\frac{M_i}{1+\exp(-4q/M_i)}-\frac{M_i}{2}=\frac{M_i}{2}\tanh\!\left(\frac{2q}{M_i}\right).
\end{aligned}
\end{equation}
Then $F_i(q)=\sigma_i(q-c_i)+M_i/2$. Define the centered states
\begin{equation*}
\tilde{x}_\mrm{S}:=\mrm{STN}-\frac{M_S}{2},
\qquad
\tilde{x}_\mrm{G}:=\mrm{GP}-\frac{M_G}{2},
\end{equation*}
and offsets
\begin{equation*}
\begin{aligned}
b_S&:=w_{CS}\mrm{Ctx}-\frac{w_{GS}M_G}{2}-c_S,\\
b_G&:=\frac{w_{SG}M_S}{2}-\frac{w_{GG}M_G}{2}
      -w_{XG}\mrm{Str}-c_G.
\end{aligned}
\end{equation*}
This yields a model in which the nonlinearities $\sigma_{\mrm{S}}$ and
$\sigma_{\mrm{G}}$ are centered and satisfy $\sigma_i(0)=0$:
\begin{equation}\label{eq:stn-gpe-centered}
\begin{aligned}
\tilde{z}_{\mrm{S}}(t)&:=-w_{\mrm{GS}}\tilde{x}_{\mrm{G}}(t-\tau_{\mrm{GS}})+b_S,\\
\tilde{z}_{\mrm{G}}(t)&:=w_{\mrm{SG}}\tilde{x}_{\mrm{S}}(t-\tau_{\mrm{SG}})
          -w_{\mrm{GG}}\tilde{x}_\mrm{G}(t-\tau_{\mrm{GG}})+b_G,\\
\tau_{\mrm{S}}\dot{\tilde{x}}_{\mrm{S}}(t)
      &=\sigma_{\mrm S}(\tilde{z}_{\mrm S}(t))
        -\tilde{x}_{\mrm{S}}(t)+\tilde{u}(t),\\
\tau_{\mrm{G}}\dot{\tilde{x}}_{\mrm{G}}(t)
      &=\sigma_{\mrm G}(\tilde{z}_{\mrm G}(t))
        -\tilde{x}_{\mrm{G}}(t).
\end{aligned}
\end{equation}
Here $\tau_{ij}$ and $w_{ij}$ denote the transmission delay and synaptic
weight, respectively, from population $i$ to population $j$.
$\tau_{\mrm{S}}$ and $\tau_{\mrm{G}}$ denote the time constants for the STN and GPe populations,
respectively. To formulate the control objective around a desired operating point, we introduce the equilibrium
 \((\tilde x_S^\star,\tilde x_G^\star)\) satisfying
\begin{align*}
      \tilde z_{\mrm{S}}^\star&:=-w_{\mrm{GS}}\tilde x_{\mrm{G}}^\star+b_S,\\
      \tilde z_{\mrm{G}}^\star&:=w_{\mrm{SG}}\tilde x_{\mrm{S}}^\star
          -w_{\mrm{GG}}\tilde x_\mrm{G}^\star+b_G,\\
      \tilde x_\mrm{S}^\star
          &=\sigma_{\mrm S}(\tilde z_{\mrm S}^\star)+\tilde u^\star,\\
      \tilde x_\mrm{G}^\star
          &=\sigma_{\mrm G}(\tilde z_{\mrm G}^\star).
\end{align*}
For $i\in\{\mrm{S},\mrm{G}\}$, introduce the deviation variables
\begin{equation*}
x_i(t):=\tilde{x}_i(t)-\tilde x_i^\star,
\quad
z_i(t):=\tilde{z}_i(t)-\tilde z_i^\star,
\quad
u(t):=\tilde{u}(t)-\tilde u^\star,
\end{equation*}
and define the shifted nonlinearities
\begin{equation*}
\delta_i(v):=\sigma_i(\tilde z_i^\star+v)-\sigma_i(\tilde z_i^\star).
\end{equation*}
The resulting system is a nonlinear DDE with the equilibrium shifted to the
origin.
Following~\cite{peet2020representationnetworkssystemsdelay}, we introduce 
the delayed states $\phi_{ij}(t,s):={x}_i(t-\tau_{ij}s)$ for
$(i,j)\in\mathcal I_{\mrm D}$ and $s\in[0,1]$, where
$\mathcal I_{\mrm D}:=\{(\mrm G,\mrm S),(\mrm S,\mrm G),(\mrm G,\mrm G)\}$.
This yields the
equivalent ODE--PDE representation,
\newlength{\philhswidth}
\settowidth{\philhswidth}{$\dot{\phi}_{ij}(t,s)\mathrel{}$}
\begin{equation}\label{eq:stn-gpe_ode-pde}
\begin{aligned}
z_{\mrm{S}}(t)
    &:= -w_{\mrm{GS}}\phi_{\mrm{GS}}(t,1),\\
z_{\mrm{G}}(t)
    &:= w_{\mrm{SG}}\phi_{\mrm{SG}}(t,1)
        -w_{\mrm{GG}}\phi_{\mrm{GG}}(t,1),\\
\tau_{\mrm{S}}\dot{x}_{\mrm{S}}(t)
    &= \delta_{\mrm S}(z_{\mrm S}(t))-x_{\mrm{S}}(t)+u(t),\\
\tau_{\mrm{G}}\dot{x}_{\mrm{G}}(t)
    &= \delta_{\mrm G}(z_{\mrm G}(t))-x_{\mrm{G}}(t),\\
&\hspace{-\philhswidth}
\left.
\begin{aligned}
\dot{\phi}_{ij}(t,s)
    &= -\frac{1}{\tau_{ij}}\phi_{ij,s}(t,s),\\
\phi_{ij}(t,0)
    &= x_i(t).
\end{aligned}
\right\}
\quad
(i,j)\in\mathcal I_{\mrm D}
\end{aligned}
\end{equation}
Thus, the equilibrium shift only affects the nonlinear operators, while leaving the linear delayed dynamics unchanged.
We now turn to robustness and performance analysis. 
\subsection{Robust Stability Analysis}
We first seek the largest sector or slope bound $\beta$ for which a
stimulation controller can certify robust stability of
\eqref{eq:stn-gpe_ode-pde}. Thus each feasibility test includes controller
synthesis: the controller is a design variable throughout the search
in $\beta$. At this stage the objective is robust stabilization, with
the performance channels omitted. Write
\begin{equation*}
z_\Delta:=\bmat{z_{\mrm S}\\z_{\mrm G}},
\qquad
w_\Delta:=\bmat{\delta_{\mrm S}(z_{\mrm S})\\
                 \delta_{\mrm G}(z_{\mrm G})}.
\end{equation*}
The shifted nonlinearities are memoryless, time invariant, monotone, and
vanish at the origin. We use a common upper bound $\beta>0$ for both
channels, first retaining only the sector information
\begin{equation}\label{eq:stn-gpe-sector-bound}
0\leq v\delta_i(v)\leq\beta v^2,
\qquad i\in\{\mrm S,\mrm G\},
\end{equation}
and then also using the slope restriction
\begin{equation}\label{eq:stn-gpe-slope-bound}
0\leq\frac{\delta_i(v_1)-\delta_i(v_2)}{v_1-v_2}\leq\beta,
\qquad v_1\neq v_2.
\end{equation}
The slope restriction implies the sector bound by taking $v_2=0$.
We initially test these uncertainty classes on the whole real line; below
we identify the trajectories of the actual sigmoids to which a given
$\beta$ applies. Here $\beta$ denotes a sector or slope bound, independently
of the beta frequency band in the example title.

\subsubsection{Sector uncertainty}
0.492 For $D:=\operatorname{diag}(d_{\mrm S},d_{\mrm G})\succ0$, the sector
IQC acting on $\bmat{z_\Delta^\top&w_\Delta^\top}^\top$ has multiplier
\begin{equation}\label{eq:stn-gpe-sector-multiplier}
\Pi_{\mrm{sec}}:=\bmat{0&\beta D\\\beta D&-2D}.
\end{equation}
Its quadratic form is nonnegative pointwise by
\eqref{eq:stn-gpe-sector-bound}. Applying Lemma~\ref{lem:sec-bound}
channelwise with $\alpha=0$ directly gives the hard dual IQC
\begin{equation*}
\underline{\Psi}_\Delta=I_4,
\qquad
\underline{\mcl V}_\Delta=\bmat{0&\beta D\\\beta D&-2D}.
\end{equation*}
Thus the sector synthesis uses this hard dual IQC directly, without an
$\epsilon$ completion or the ARE-based factorization. For each trial
$\beta$ and an admissible choice of $D$, we seek a controller and storage
operator satisfying the stability-only synthesis LPI with this IQC.
Searching in $\beta$ gives the
largest sector bound certified by this design procedure, denoted
$\beta_{\mrm{sec}}^*$, together with a stabilizing controller at each
feasible bound.

\subsubsection{Zames--Falb multiplier}
Zames--Falb multipliers exploit monotonicity and slope restrictions on
nonlinearities~\cite{zames1968stability}.
To exploit~\eqref{eq:stn-gpe-slope-bound}, 0.562 replace the static weight $D$
by the diagonal Zames--Falb (ZF) multiplier, here we use the construction provided in \cite{turner2022zamesfalbmultipliersdontpanic},
\begin{equation}\label{eq:stn-gpe-zf-transfer}
\begin{aligned}
M(\zeta)&:=I-\operatorname{diag}(H_{\mrm S}(\zeta),H_{\mrm G}(\zeta)),\\
H_i(\zeta)&:=\sum_{k=1}^{3}\frac{\kappa_{ik}}{\zeta+a_{ik}},\\
a_{ik}&>0,\qquad\kappa_{ik}\geq0,\qquad
\sum_{k=1}^{3}\frac{\kappa_{ik}}{a_{ik}}\leq1.
\end{aligned}
\end{equation}
Each term has a nonnegative exponential impulse response, with integral
weight $\kappa_{ik}/a_{ik}$. Thus $\|H_i\|_{L_1}\leq 1$, giving a scaled ZF multiplier
valid for monotone, slope-restricted nonlinearities without assuming
oddness~\cite[Thm.~4]{turner2022zamesfalbmultipliersdontpanic}.
This matters because shifting the sigmoid to an arbitrary equilibrium
generally removes its odd symmetry.

We choose the three pole rates from the population time constants and
the recurrent GPe delay:
\begin{equation}\label{eq:stn-gpe-zf-coefficients}
\bmat{a_{i1}&a_{i2}&a_{i3}}
=c_i\bmat{\tau_{\mrm G}^{-1}&\tau_{\mrm S}^{-1}&\tau_{\mrm{GG}}^{-1}},
\qquad c_i>0.
\end{equation}
The poles $-a_{ik}$ supply memory on the GPe relaxation, STN relaxation,
and GPe recurrent-delay scales. The reciprocal delay is a characteristic
rate, not a pole of the delay operator. The dimensionless factors $c_i$
adjust each multiplier's time scale relative to these physical scales;
we initially take $c_{\mrm S}=c_{\mrm G}=1$.

To relate the residues to the synaptic coupling strengths, define
\begin{equation*}
\begin{aligned}
\bmat{h_1&h_2&h_3}
&:=\frac{\rho\bmat{|w_{\mrm{GS}}|&|w_{\mrm{SG}}|&|w_{\mrm{GG}}|}}
{|w_{\mrm{GS}}|+|w_{\mrm{SG}}|+|w_{\mrm{GG}}|},\\
\kappa_{ik}&:=a_{ik}h_k,\qquad 0<\rho<1.
\end{aligned}
\end{equation*}
This assigns the GPe relaxation term a weight proportional to the
GPe-to-STN coupling, the STN relaxation term a weight proportional to
the STN-to-GPe coupling, and the recurrent-delay term a weight
proportional to the GPe self-coupling. The weights are nonnegative,
and normalization preserves the ZF norm condition as the synaptic
parameters change:
\begin{equation*}
\sum_{k=1}^{3}\frac{\kappa_{ik}}{a_{ik}}
=\sum_{k=1}^{3}h_k=\rho<1,
\qquad\|H_i\|_{L_1}=\rho.
\end{equation*}
We take $\rho=13/15$ to retain the prescribed kernel norm. This
parameterization is a physically motivated choice of multiplier shape;
it does not assert that the synaptic strengths determine the optimal
IQC weights or guarantee an improved certificate. The factors $c_i$
and $\rho$ can be compared through an outer synthesis search. For
fixed system parameters and these choices, the poles and residues
remain fixed during the search in $\beta$. The resulting slope IQC is
\begin{equation}\label{eq:stn-gpe-zf-multiplier}
\Pi_{\mrm{ZF}}(j\omega):=\bmat{
0&\beta M(j\omega)^*\\
\beta M(j\omega)&-\bigl(M(j\omega)+M(j\omega)^*\bigr)}.
\end{equation}
With all $\kappa_{ik}=0$, this reduces to the sector multiplier. We repeat
the search in $\beta$, synthesizing a controller for each trial bound
using the ZF family, and denote the largest certified bound by
$\beta_{\mrm{ZF}}^*$. Retaining the sector solution as a candidate gives
$\beta_{\mrm{ZF}}^*\geq\beta_{\mrm{sec}}^*$ when the controller and
storage parameterizations include those of the sector design. These
bounds describe what the chosen synthesis procedure certifies; failure
of a synthesis LPI does not establish that robust stabilization is
impossible.

For the ZF multiplier, we use a $J$-spectral factorization to obtain
the required hard IQC and dual filter. Since the sector bound gives
$\|w_\Delta\|\leq\beta\|z_\Delta\|$, we use the PN completion as discussed in \cite{6915700}
\begin{equation}\label{eq:stn-gpe-combined-multiplier}
\Pi_\Delta:=\Pi_{\mrm{ZF}}+
\epsilon\bmat{\beta^2I_2&0\\0&-I_2},\qquad\epsilon>0,
\end{equation}
The added term is a valid norm IQC. Since
$M(j\omega)+M(j\omega)^*\succeq0$, $\Pi_\Delta$ is strict PN and
Theorem~\ref{thm:strict-pn-j-spectral} supplies
\begin{equation*}
\Pi_\Delta=\Theta_\Delta^\sim J_{2,2}\Theta_\Delta,
\qquad
\Theta_\Delta,\Theta_\Delta^{-1}\in\mathbb{RH}_\infty.
\end{equation*}
For each trial $\beta$ in the ZF search, select admissible multiplier
weights and compute this factorization. With $\Theta_\Delta$ fixed, the dual synthesis LPI
in Corollary~\ref{thm:state-feedback}, with performance channels omitted,
searches for a controller and storage operator certifying closed-loop
robust stability. Multiplier choices can be compared outside this LPI;
joint optimization of the factorization and controller is not assumed
to be convex. The resulting feasible values of $\beta$ therefore belong
to controller--certificate pairs.

\subsubsection{Admissible trajectories}
For a controller certified at a given $\beta$, the admissible sigmoid
inputs follow from~\eqref{eq:stn-gpe-centered-sigmoid}:
\begin{equation*}
\delta_i'(v)={\text{sech}^2\!\left(2(\tilde z_i^\star+v)/M_i\right)}.
\end{equation*}
Consequently the sector and slope bounds hold globally for $\beta\geq1$.
For $0<\beta<1$, inversion of this derivative gives
\begin{equation*}
\delta_i'(v)\leq\beta
\quad\Longleftrightarrow\quad
|\tilde z_i^\star+v|\geq
\frac{M_i}{2}\operatorname{arcosh}(\beta^{-1/2}).
\end{equation*}
To obtain an interval containing the equilibrium, remain on its side of
the sigmoid midpoint. A simple sufficient trajectory bound is
\begin{equation}\label{eq:stn-gpe-admissible-trajectories}
|z_i(t)|<|\tilde z_i^\star|
          -\frac{M_i}{2}\operatorname{arcosh}(\beta^{-1/2}),
\end{equation}
for $i\in\{\mrm S,\mrm G\}$ and every $t\geq0$.
The right-hand side is positive when $\beta>\delta_i'(0)$. Throughout
this interval the derivative is bounded by $\beta$, so both
\eqref{eq:stn-gpe-sector-bound} and~\eqref{eq:stn-gpe-slope-bound} hold.
For an equilibrium below the midpoint, the less restrictive one-sided
condition is
\begin{equation*}
z_i(t)\leq-\tilde z_i^\star-
\frac{M_i}{2}\operatorname{arcosh}(\beta^{-1/2}).
\end{equation*}
Above the midpoint, the corresponding condition is
\begin{equation*}
z_i(t)\geq-\tilde z_i^\star+
\frac{M_i}{2}\operatorname{arcosh}(\beta^{-1/2}).
\end{equation*}
These bounds constrain the delayed firing-rate combinations
in~\eqref{eq:stn-gpe_ode-pde}. A larger feasible $\beta$ permits a wider
interval around the same equilibrium. The local IQC certificate applies
via a globally slope-restricted extension agreeing with the sigmoid on
this interval, and transfers to the original system for trajectories
that remain there; invariance requires a separate check.

\subsection{Robust Performance Analysis}
Having established the robustness limit through controller synthesis,
fix a feasible $\beta<\beta_{\mrm{ZF}}^*$ and synthesize controllers
with performance objectives at this smaller bound. This leaves margin
relative to the robust-stability limit. For a local sigmoid
certificate, also require $\beta>\max_i\delta_i'(0)$ so that
\eqref{eq:stn-gpe-admissible-trajectories} contains a neighborhood of the
equilibrium. This choice is possible only if the two requirements leave
a nonempty interval. Keep $\beta$ fixed throughout the performance
comparison so that all controllers are assessed on the same admissible
trajectory range.

Let $w_{\mrm p}$ be an additive disturbance in the stimulation channel.
Using the same uncertainty channels and transport equations as above,
the generalized plant is specified by
\begin{equation}\label{eq:stn-gpe-generalized-plant}
\begin{aligned}
\tau_{\mrm S}\dot{x}_{\mrm S}
    &=w_{\Delta,1}-x_{\mrm S}+u+w_{\mrm p},\\
\tau_{\mrm G}\dot{x}_{\mrm G}
    &=w_{\Delta,2}-x_{\mrm G},\\
z_{\mrm p}&:=\bmat{x_{\mrm S}\\x_{\mrm G}\\\mu u},\qquad\mu>0.
\end{aligned}
\end{equation}
Thus the performance objective measures rejection of perturbations
entering through the stimulation input. The effort penalty concerns
the stimulation deviation $u=\tilde u-\tilde u^\star$.
We assess the induced $L_2$ gain to each signal separately, using
$\gamma_{\mrm S}$, $\gamma_{\mrm G}$, and $\gamma_u$ for
$x_{\mrm S}$, $x_{\mrm G}$, and $u$, respectively.
To generate controllers with different effort penalties, for zero
initial deviations and controller state minimize the weighted synthesis
bound $\gamma$ in
\begin{equation}\label{eq:stn-gpe-performance-bound}
\|x_{\mrm S}\|_{L_2}^2+\|x_{\mrm G}\|_{L_2}^2
+\mu^2\|u\|_{L_2}^2\leq\gamma^2\|w_{\mrm p}\|_{L_2}^2.
\end{equation}
This directly penalizes firing-rate deviations and control effort,
with the performance IQC
\begin{equation}\label{eq:stn-gpe-performance-filter}
\Psi_{\mrm p}=I_4,\qquad
\mcl V_{\mrm p}=\operatorname{diag}(I_3,-\gamma^2).
\end{equation}
For $\gamma>0$, its $J_{3,1}$-spectral factor
$\operatorname{diag}(I_3,\gamma)$ is stably invertible and can be combined
with $\Theta_\Delta$ in the dual synthesis condition.

For each $\mu$, fix an admissible uncertainty multiplier at the chosen
$\beta$, compute its factorization, and use
Corollary~\ref{thm:state-feedback} to synthesize the stimulation controller
while minimizing $\gamma$. Keeping the multiplier fixed in each synthesis
problem avoids a joint search over its factorization and the controller.
Increasing $\mu$ gives more weight to reducing stimulation effort;
decreasing it gives more weight to suppressing $x_{\mrm S}$ and
$x_{\mrm G}$. We will vary $\mu$ to obtain a family of controllers and
then reanalyze each fixed controller with three separate performance
outputs to certify
\begin{equation*}
\begin{aligned}
\|x_{\mrm S}\|_{L_2}&\leq\gamma_{\mrm S}\|w_{\mrm p}\|_{L_2},\\
\|x_{\mrm G}\|_{L_2}&\leq\gamma_{\mrm G}\|w_{\mrm p}\|_{L_2},\\
\|u\|_{L_2}&\leq\gamma_u\|w_{\mrm p}\|_{L_2}.
\end{aligned}
\end{equation*}
Each gain bound is minimized in its own analysis problem, for the same
controller, disturbance channel, and uncertainty bound $\beta$.
We will report all three bounds and compare $\gamma_u$ with
$\gamma_{\mrm S}$ and $\gamma_{\mrm G}$ separately, showing how the
effort--attenuation tradeoff differs between the two populations.
The weighted $\gamma$ serves only as the synthesis objective; the
reported gains refer to the individual, unweighted signals.
For the original sigmoid with
$\beta<1$, each comparison is restricted to trajectories satisfying
\eqref{eq:stn-gpe-admissible-trajectories}; the $L_2$-gain bounds alone
do not establish this pointwise condition.
